% Nouveau EGA V (cote 25): archive cote 25, site page 89 (PDF page 90).
% typescript. Draft compared: none.
% AI-transcribed by Claude Opus 5.5 (claude-opus-5-5, adaptive thinking, effort high), Message Batches API, batch ega25b (msgbatch_01A7n5YuXubUPNK1zsepi8RG), 2026-11-13; not yet checked by eye.
\scanpage{25/89}
% NOTE: the top right corner is cut off at the edge of the scan; only the lower part of a number is visible there (probably the page number)
seule conclusion $p_C$ plat, l'hypothèse CM sur X par une simple hypothèse
dimensionelle. Supposons par exemple f une immersion et X irr de dim m,
de sorte que $p_C$ est qu fini , et comme $X_{k\hand{(y)}}$ et $Q(C_{\hand{y}})$ sont irréd de
même dimension, \xout{le premier CM} et le deuxième régulier, $p_C$ ne peut être
plat que si $X_{k\hand{(y)}}$ est CM.

Des propriétés plus fines sont les propriétés différentielles de
la projection conique, notamment pour X lisse sur k et $f:X\longrightarrow P$ non rami-%
\inserted[typed]{étudiées au nº12.}
fié, \xout{Dans le cas où les composantes} Rappelons que \xout{la fibre générique de}
en dehors d'une partie Z de codimension 1 de $Q(C\ )$, le morphisme
$p_{C_{\hand{y}}}$ \inserted[typed]{sur $X(C_{\hand{y}})$} est \uline{lisse}, Et une analyse plus détaillée (esquissée au nº12) montre
que si \uline{les dimensions des composantes} de X sont $\geq$ m, alors\hand{)}
(ou montrerait, si nous ne la faisons pas) \xout{que} en dehors d'une partie $Z'\subset Z$
de $Q(C_{\hand{y}}\ )$ de codimension $\geq 2$, les fibres $p_C^{-1}(\hand{\zeta}) = \underline{X}^{(m)}_{\hand{\zeta}}$ ne peuvent
avoir au pis que des points \xout{doub} singuliers ordinaires (au s ens
géométrique), et en fait (si f est une immersion et X est géom irréd),
\xout{juste} au plus \uline{un} tel point, ce dernier \inserted[typed]{étant alors} \xout{\unsure{ayant}} nécessairement \xout{\illegible{1}} .
rationnel sur $k(\hand{\zeta})$, - ces assertions étant valables tout au moins si
k est de car 0, \uline{ou} à condition\xout{s} de remplacer f par un $\phi_n f$ $(n\geq 2)$ comme
au nº 9. Il y a lieu également de donner des propriétés différentielles
de $p_{C_{\hand{y}}}$ dans le cas où dim X $\leqslant$ m , et où par suite $p_{C_{\hand{y}}}$ est défini sur $X_{k\hand{(y)}}$.
Je me borne à indiquer les propriétés suivantes, dont la démonstration
\inserted{doit} être facile et \inserted{est} laissée à Dieudonné:

\uline{Proposition} 14.4. Supposons $f:X\longrightarrow P$ non ramifié, et dim X $\leqslant$ m.
Soit T un sous-préschéma fini de X, Alors :\xout{la restr}
\margin{\unsure{X?}} a) \inserted{Si f est une immersion,} La restriction de $p_C$ à $T_{k\hand{(y)}}$ est \deleted{un monomorphisme,} \inserted{radiciel i.e.\ ``géométriquement injectif''.} Si de
plus Y est une partie fermée de X de dimension $\leqslant$ m-1, on a
$p_{C_{\hand{y}}}^{-1}(p_{C_{\hand{y}}}(Y_{k\hand{(y)}})) \cap T_{k\hand{(y)}} = \emptyset$ .
b) Si X est lisse en les points de T, alors $p_{C_{\hand{y}}}$ est non ramifié
en tous point \xout{de} $T_{k\hand{(y)}}$, \xout{au dessus de tout point de T} \inserted{et \unsure{même|aussi} en les pts de $p_{C_y}^{-1}\,p_{C_y}(T_{k(y)})$.}
\mspage{\hand{83}}
